% Page 009

\seite{9}

\margmark{Ritter von d. Els,\\Else, von der Leyen\\Leyen.\\Urweiche.}%
aus dem zwölften Jahrhundert, laut welcher ein Ritter\ob
von der Els und seine Ehefrau Else von der\ob
Leyen beider zu Urweiche liegenden Ländereien,\ob
Wiesen etc der Pfarrkirche zu\unsicher{} Seffern g. schenkten\ob
unter der Bedingung, daß die Güter im Besitze der\ob
Geschenkgeber bis zu ihrem Tode verbleiben sollte,\ob
bis dahin aber verpflichteten sich die Geschenkgeber,\ob
der Kirche jährlich zwei Trierer Malter Hafer (altes\ob
Maß etwa fünfzehn Scheffel) zu liefern. Sie nahmen\ob
also ihre eigenen Güter von der Kirche zu Lehen. Größe\ob
und Umfang dieser Güter sind nicht angegeben, doch\ob
müssen dieselben sehr bedeutend \darueber{enthielt 5 M.} gewesen sein. Aber\ob
wer waren die Geschenkgeber und wie wurde ein\ob
Ritter von der Els und eine Else von der Leyen [darüber korrigiert: Leyen]\ob
nach Urweiche verschlagen. Ein Schloß Els besteht\ob
heute noch im Regierungsbezirk Koblenz an dem\ob
Flüßchen Elz. Das Geschlecht der Leyen aber muß\ob
schon längst den Weg alles Fleisches gegangen sein.\ob
Keine Spur ist heute hier mehr zu finden. Nur aus\ob
dem Munde alter Leute, die es ebenfalls nur aus\ob
der Tradition wissen können, hat man vernom\ohb
men, daß auf der „Heide“ südöstl. [darüber eingefügt] von S.\ ein Schloß\ob
gestanden haben soll. Wer weiß es? Wer hat's ge\ohb
sehen? Seine Herrlichkeit ist untergegangen und\ob
der Wind pfeift nicht mehr durch seine Hallen. Der\ob
Pflug hat schon längst seine Furchen hinübergezogen,
\pend
